\section{How to Read and Reproduce the Current Method}
\label{app:causal-walkthrough}

This companion begins with the recovery-aware causal method used in the
current article. Sections~\ref{app:causal-walkthrough}--\ref{app:current-scenarios}
give its execution order, local derivations, worked examples and simulation
contracts. The separately compiled Historical Technical Notes preserve
full diagnosability proofs, earlier state-reconstruction and batch-estimator
studies, and secondary evidence. The earlier studies support particular
components; they are not additional tests of the complete causal algorithm.\AIDraftingCredit{}

The article retains the physical model, the defining evidence integral,
the recovery mechanism, the decision outputs and decisive results. This
supplement supplies the intermediate calculations. The Historical Technical
Notes develop functional observability, nuisance-profiled information and
source ambiguity, and distinguish the implemented batch estimator $\mathcal E_{\mathrm B}$ from
the proposed joint estimator $\mathcal E_{\mathrm J}$. Neither should be
silently substituted for the causal history bank, archived as R11.
The conditional Gaussian updates use standard filtering identities
\cite{sarkka2013bayesian}; the derivation below specifies their role and
the approximations made by the implemented bank.

\subsection{Inputs, Outputs and the Information Boundary}

The inputs are successive reports from eight PMUs, the fixed physical and
instrumental operators, the admitted load/line catalog and the frozen
method configuration. The output is a timestamped sequence of detection,
family/source and restoration messages, together with retained attribution
masses. True ON/OFF times, true sources and generating severities belong
to the evaluator. They are not inputs to the causal algorithm.

Three times must remain distinct: the latest observed report, the estimated
physical onset, and the public emission time. Looking back to propose an
onset is compatible with causality when the message is emitted only after
the required reports arrive. It does not permit backdating the public alarm.

\subsection{Execution Order for One Stream}

\begin{enumerate}
\item \textbf{Load the frozen operators.} Retain voltage reference,
oriented terminal currents, load input directions, switch modes and
instrument memory. A line opening changes the network and observation
operators. It is not represented by fitting a load amplitude.
\item \textbf{Fit the startup context.} Use the first 300 reports for
operating factors, common reference and noise scale. For each network
hypothesis, set the voltage/current discrepancy to the difference between
the observed prefix mean and its reference prediction. Keep that offset
fixed afterwards. Frequency and ROCOF use their empirical prefix centers.
The prefix is not scored again as independent model-selection evidence.
\item \textbf{Predict the next report.} Propagate the retained histories
and evaluate each network model's predictive density before correcting
states or choosing new histories with the current report.
\item \textbf{Propose candidate onsets.} The frozen screen uses lags
0--18 reports within a ten-second stored state history. Each proposal
starts from a stored state before its proposed onset, including any
recovery already present. The bank does not enumerate every possible
onset in the stored interval.
\item \textbf{Score histories.} For loads, integrate severity with the
bounded prior and the local conditional Gaussian likelihood. For lines,
score the admitted switched operator. Score a report before using that
same report for its conditional state correction.
\item \textbf{Update and compress.} Update severity and conditional
states, network weights, restoration children and the retained beam.
Record discarded mass. At episode closure retain the augmented-state
mixture mean; the omitted between-history spread is a documented
approximation, not an exact posterior covariance.
\item \textbf{Publish and record.} Aggregate active mass $M$, cause
mass $\Gamma$ and $\Delta M$. Apply the frozen detection, diagnosis and
restoration rules. Save emission time and the last report used, in addition
to estimated onset and source. Continue propagation through OFF.
\end{enumerate}

\subsection{Why Replayed Reports Are Not New Evidence}

Let a retained parent have log score
$\ell_{0:k}=\ell_{0:\tau-1}+\ell^{\rm old}_{\tau:k}$.
Replacing its suffix by a newly proposed physical history gives
\begin{equation}
 \ell^{\rm new}_{0:k}
 =\ell_{0:k}-\ell^{\rm old}_{\tau:k}+\ell^{\rm new}_{\tau:k}.
 \label{eq:supp-suffix-replacement}
\end{equation}
Transition and prior terms must accompany their corresponding history.
The rolling screen supplies a proposal, not an extra independent
likelihood. Adding the new suffix without removing the old suffix would
count the same reports twice. This bookkeeping statement does not make
adaptive screening or beam pruning exact Bayesian inference.

\subsection{Frozen Numerical Choices}

The input grid is 30 reports/s. The five network hypotheses are nominal,
global $\pm2\%$ and selected-branch $\pm3\%$ paired R/X perturbations.
Each hypothesis is fixed in time; the weights change. They do not identify
a continuous impedance posterior.

The per-report birth and restoration hazards are
$1-\exp[-(1/60)/30]$ and $1-\exp[-(1/30)/30]$.
Birth mass is divided equally between load and line families, then
uniformly among the admitted sources within each family. Load severity
has a zero-mean Gaussian prior with standard deviation 0.15, truncated to
$[-0.2,0.2]$ and integrated over intervals bounded by
$-0.2,-0.1,0,0.1,0.2$. After three reports, eligible active histories can
create OFF children at either half-report position, inheriting state and
severity. After likelihood normalization, the beam retains the union of
the 96 highest-weight recent histories, 64 older histories and at least
one per occupied source. The union is renormalized.

Active mass 0.99 for six reports confirms detection; restoration mass
0.99 for 15 reports confirms OFF. Source/line confirmation and episode
replacement additionally follow the archived publication configuration.
These thresholds are implementation settings, not guarantees of
probability calibration or false-alarm control. Nominal and mixture use
their frozen publication paths, so their comparison does not isolate
network averaging alone. The admission-rule ablation changes publication,
not the recovery dynamics.

\section{Local Calculations and Worked Examples}
\label{app:causal-calculations}

\subsection{Completing the Bounded-Amplitude Integral}
\label{sec:supp-causal-integral}

For one retained severity interval, write
$e=sa+\epsilon$, with $\epsilon\sim\mathcal N(0,S)$,
$S\succ0$, and $a\sim\operatorname{TN}_{[l,u]}(\mu,v)$, $v>0$.
Here $\mu,v$ are parameters of the underlying untruncated Gaussian, not
the bounded distribution's moments. With
$g=s^\top S^{-1}s$ and $q=s^\top S^{-1}e$, expansion gives
\begin{align}
 &(e-sa)^\top S^{-1}(e-sa)+(a-\mu)^2/v\nonumber\\
 &\quad=e^\top S^{-1}e+\mu^2/v-\mu'^2/v'
                  +(a-\mu')^2/v',\\
 &v'=(v^{-1}+g)^{-1},\qquad
 \mu'=v'(\mu/v+q).
\end{align}
Define $Z(\mu,v)=\Phi((u-\mu)/\sqrt v)-\Phi((l-\mu)/\sqrt v)$.
Integrating the last Gaussian square over $[l,u]$ yields
\begin{align}
 \log L={}&-\tfrac12[e^\top S^{-1}e+\log|S|+n_y\log(2\pi)]\nonumber\\
 &+\tfrac12[\log(v'/v)+\mu'^2/v'-\mu^2/v]\nonumber\\
 &+\log Z(\mu',v')-\log Z(\mu,v).
 \label{eq:supp-causal-integral}
\end{align}
The posterior on that interval is
$\operatorname{TN}_{[l,u]}(\mu',v')$. With
$\alpha=(l-\mu')/\sqrt{v'}$, $\beta=(u-\mu')/\sqrt{v'}$,
$Z'=\Phi(\beta)-\Phi(\alpha)$ and standard normal density $\phi$,
its mean is
\begin{equation}
 \mathbb E[a\mid e]=\mu'+\sqrt{v'}\,
       \frac{\phi(\alpha)-\phi(\beta)}{Z'}.
\end{equation}
Thus neither the untruncated mean $\mu'$ nor a best-fit severity can simply
replace the integral. The implementation combines retained intervals using
their weights. The identity is exact for this conditional Gaussian model;
it does not remove linearization, shared-covariance or pruning error.

\subsection{Numerical Example of One Interval}

This dimensionless example illustrates the calculation; it is not an
IEEE-39 simulation or a new performance result. Set
$S=1$, $s=2$, $e=1$, $\mu=0$, $v=1$ and $[l,u]=[0,1]$.
Then $g=4$, $q=2$, $v'=0.2$ and $\mu'=0.4$. Substitution in
\eqref{eq:supp-causal-integral} gives
$L\simeq0.342688$ and
$\mathbb E[a\mid e]\simeq0.464941$.
The different posterior mean follows from truncation.
The script \nolinkurl{scripts/verify_supplement_examples.py} checks the
closed form against direct numerical integration, including the posterior
mean. It checks algebra only and does not modify a frozen experiment.\AIDraftingCredit{}

\subsection{Conditional State Correction}
\label{sec:supp-causal-state}

Represent the conditional state mean as $m_h+d_h a$. Let
$T_o,J_o,g_o,u_o,j_o,v_o$ denote the archived one-report transition and
observation operators for switch $o$. The code forms
$e=r_k-J_o m_h-j_o$ and $s=J_o d_h+v_o$. With shared covariance
$P_\rho$ in discrepancy regime $\rho$,
\begin{align}
 S_o&=R+J_oP_\rho J_o^\top,&
 K_o&=T_oP_\rho J_o^\top S_o^{-1},\nonumber\\
 m_h^+&=T_o m_h+g_o+K_o e,&
 d_h^+&=T_o d_h+u_o-K_o s.
 \label{eq:implemented-conditional-correction}
\end{align}
To see the last line, apply the correction to a fixed amplitude:
\[
 T_o(m_h+d_h a)+g_o+u_o a+K_o(e-sa)
     =m_h^++d_h^+a.
\]
The severity distribution is updated by the same likelihood integral.
Regime covariances come from nominal-operator Riccati equations. R11 does
not propagate an independent exact covariance for every switched history.
At closure the mean is retained and the omitted mixture spread recorded.

\subsection{Restoration Is Not a State Reset}

Consider the illustrative scalar system
$\dot x=-x+u$, $y=x$, with $x(0)=0$. Apply $u=a$ for
$0\le t<\kappa$ and restore $u=0$ afterwards. Direct integration gives
\begin{equation}
 x(t)=\begin{cases}
 a(1-e^{-t}),&0\le t<\kappa,\\
 a(1-e^{-\kappa})e^{-(t-\kappa)},&t\ge\kappa.
 \end{cases}
\end{equation}
Both expressions agree at $\kappa$. For $a=1$ and $\kappa=2$,
$x(2)=0.864665$ and $x(3)=0.318092$ even though the intervention is OFF.
Resetting the reference to equilibrium would discard that predictable
recovery. This example explains state continuity; it does not claim that
one exponential describes the multimode IEEE-39 response.

\subsection{Active Mass, Cause Mass and Model Reweighting}

For one model, suppose an illustrative cell has mass 0.60 in active
histories and 0.30 in restored histories, with 0.10 elsewhere. Its active
mass is $M_c=0.60$, while its originating-cause mass is $\Gamma_c=0.90$.
Recovery can therefore remain attributable without declaring the input
active. These assigned masses illustrate definitions, not calibrated
probabilities from an experiment.

For two reports, adding and subtracting
$\sum_\ell w_{\ell,k}M_{\ell,c,k-1}$ proves
\begin{align}
 \Delta M_c={}&\sum_\ell w_{\ell,k}
       (M_{\ell,c,k}-M_{\ell,c,k-1})\nonumber\\
 &+\sum_\ell(w_{\ell,k}-w_{\ell,k-1})M_{\ell,c,k-1}.
\end{align}
For example, unchanged conditional masses $(0.8,0.2)$ and network weights
changing from $(0.5,0.5)$ to $(0.75,0.25)$ move $M_c$ from 0.50 to 0.65.
The entire 0.15 change comes from model reweighting. A rising aggregate
mass need not mean that each model found stronger within-model evidence.

\section{Simulation Scenarios and Reproduction Order}
\label{app:current-scenarios}

\subsection{Plant, Measurement and Provenance}

The active simulator is the frozen research candidate
\texttt{SIM0001\_JuliaPD\_F01}, not an officially validated equivalent of
RAW0001. Its contract fixes Julia 1.11.9, PowerDynamics 5.0.0,
NetworkDynamics 1.3.0 and IEEE39 0.1.0. Initial states, effective parameter
vectors, terminal map, reporting operator and stochastic model come from
that snapshot. Nominal inventory files do not replace the effective
operating point. The load law is the frozen impedance law.\AIDraftingCredit{}

The voltage buses are $\{2,5,6,10,19,22,29,39\}$. Each current is an
oriented branch-terminal measurement, not a nodal injection. BUS5 retains
the provisional historical 5--6 branch, terminal 5, hypothesis; producer
CT identity is unconfirmed and BUS39 remains partially validated.
Execution uses full-precision configuration files rather than rounded
tables. RAW0001 provided retrospective calibration context. RAW0002
remains outside these studies.

Each physical parent is integrated continuously through its interventions.
Differential states remain continuous at switching, while algebraic values
are recomputed consistently. Reports use $t=k/30$ with endpoints included;
the switch-aligned report uses the post-switch algebraic value. Solver
steps are recorded separately and observations are not interpolated across
a switch. The causal instrumental history is preserved through OFF.

\subsection{Continuous Development and Confirmation}

H03 and H04 each contain three 80-s event parents with six interventions
and one 80-s normal parent. Crossing the four parents with three R/X
conditions gives 12 physical trajectories, 18 unique interventions and
54 event realizations per campaign. Both are development evidence in
the current article. Repeated conditions are not independent event draws.

C01 prespecifies 18 continuous parents, 27 persistent-line parents
(nine lines, three repetitions), and 12 normal parents. The three R/X
conditions produce $3(18+27+12)=171$ intended physical trajectories;
two instrumental-noise scales yield 342 observed records. The continuous
parents contain $18\times6=108$ unique interventions. Their six
condition/noise realizations are repeated observations of those parents.
The 120-s and 600-s grids contain 3601 and 18001 reports, respectively.
\IfFileExists{generated/c01_macros.tex}{Completed results appear in the
dedicated confirmation tables.}{At this document snapshot, complete
confirmation artifacts have not been imported; these counts specify the
intended design and are not completed performance results.}

Continuous evaluation pairs an alarm temporally before inspecting its
predicted label. The first unmatched post-ON alarm is eligible up to the
next ON, with the inherited 30-s cap. Timely attribution requires correct
family and asset emitted before physical OFF. Unmatched alarms, false
restorations and open episodes are retained. Persistent-line evaluation
instead scores the first stopping decision, including early stops,
wrong sources, omissions and censoring. Chen2016 is eligible on that lane;
it has no load or restoration output. Confidence intervals use 10,000
paired parent-cluster resamples, not independent PMU samples.

\subsection{Separate Signature and Recovery Controls}

The E\_i/X02 signature study uses a single load at buses
$\{3,7,16,24\}$, a 10-s record, known ON at 3 s and OFF at 6 s,
and one scalar operating-context direction. It contains 57 dictionary/check
trajectories and 85 test trajectories (80 interventions and five controls),
with three common noise seeds producing 255 observed test records.
The context and severity are estimated from measurements, while onset
and family are supplied. The generating source is used to assess the
response approximation and severity error. Those conditional metrics do
not measure unknown-onset causal attribution.

The X04 recovery geometry control uses 38 trajectories. Prior load
histories at buses 7 and 24, followed by delays of 0.5 and 2 s,
give four recovery contexts. Three source-pair comparisons within each
context give 12 correlated comparisons. The prior input is active from
2 to 3 s; the subsequent response is examined over 1.5 s. The supplied
history, onset, family and two-candidate universe define an oracle control,
with one uncertain prior-severity direction. It does not evaluate R11,
unknown onset, event cardinality or operational false alarms.

\subsection{Reproduction Without Changing the Experiment}

\begin{enumerate}
\item Read each campaign's frozen plan and manifest. Select the intended
population before interpreting any denominator. Keep current causal,
signature, recovery-control and historical batch studies separate.
\item Resolve the F01 state, parameter and measurement files from its
snapshot. For any new simulation create a new run directory; a new
hypothesis requires a successor candidate. Do not run historical preparation
scripts inside a frozen baseline.
\item Preserve ON/OFF continuity, report endpoints, noise seeds and the
boundary between observations and reserved event truth. Run inference
before releasing truth to the evaluator.
\item Retain solver/inference failures in intended denominators. Aggregate
at the declared parent level and preserve the original comparators and
publication rules. Do not retune from confirmation results.
\item Rebuild the two documents from the versioned tables and figures with
\nolinkurl{scripts/build_manuscript_bundle.ps1}. Its default compilation
does not regenerate scientific artifacts or import new campaign results.
An explicit artifact refresh is a separate operation.
\end{enumerate}

The repository keeps the current walkthrough in this appendix, the frozen
campaign evidence under \texttt{generated/frozen/}, and the earlier
mathematical development in the separately compiled Historical Technical Notes. A successful PDF
build verifies document consistency; it is not a new scientific validation.
\FloatBarrier
